\documentclass[11pt,a4paper]{article}

% -------------------------------------------------
% Packages
% -------------------------------------------------
\usepackage[margin=2.2cm]{geometry}
\usepackage{kotex}
\usepackage{amsmath,amssymb}
\usepackage{booktabs}
\usepackage{array}
\usepackage{xcolor}
\usepackage{listings}
\usepackage[most]{tcolorbox}
\usepackage{enumitem}
\usepackage{hyperref}
\usepackage{fancyhdr}
\usepackage{titlesec}

% -------------------------------------------------
% Colors
% -------------------------------------------------
\definecolor{codebg}{RGB}{247,247,247}
\definecolor{codeframe}{RGB}{210,210,210}
\definecolor{keywordcolor}{RGB}{0,70,160}
\definecolor{commentcolor}{RGB}{0,120,80}
\definecolor{stringcolor}{RGB}{170,50,50}
\definecolor{titleblue}{RGB}{35,75,130}

% -------------------------------------------------
% Hyperlinks
% -------------------------------------------------
\hypersetup{
    colorlinks=true,
    linkcolor=titleblue,
    urlcolor=titleblue
}

% -------------------------------------------------
% Header / Footer
% -------------------------------------------------
\pagestyle{fancy}
\fancyhf{}
\lhead{Python Programming}
\rhead{Day 7}
\cfoot{\thepage}
\setlength{\headheight}{14pt}

% -------------------------------------------------
% Section style
% -------------------------------------------------
\titleformat{\section}
  {\Large\bfseries\color{titleblue}}
  {\thesection.}{0.6em}{}

\titleformat{\subsection}
  {\large\bfseries}
  {\thesubsection.}{0.5em}{}

% -------------------------------------------------
% Python code style
% -------------------------------------------------
\lstdefinestyle{pythonstyle}{
    language=Python,
    backgroundcolor=\color{codebg},
    frame=single,
    rulecolor=\color{codeframe},
    basicstyle=\ttfamily\small,
    keywordstyle=\color{keywordcolor}\bfseries,
    commentstyle=\color{commentcolor},
    stringstyle=\color{stringcolor},
    numbers=left,
    numberstyle=\tiny\color{gray},
    numbersep=8pt,
    tabsize=4,
    showstringspaces=false,
    breaklines=true,
    columns=fullflexible,
    keepspaces=true
}

\lstset{style=pythonstyle}

% -------------------------------------------------
% Boxes
% -------------------------------------------------
\newtcolorbox{learningbox}{
    colback=blue!4,
    colframe=titleblue,
    title=\textbf{학습 목표},
    fonttitle=\bfseries
}

\newtcolorbox{importantbox}{
    colback=yellow!8,
    colframe=orange!70!black,
    title=\textbf{핵심 포인트},
    fonttitle=\bfseries
}

\newtcolorbox{practicebox}{
    colback=red!3,
    colframe=red!55!black,
    title=\textbf{실습},
    fonttitle=\bfseries
}

% -------------------------------------------------
% Document
% -------------------------------------------------
\begin{document}

\begin{titlepage}
    \centering
    \vspace*{3cm}

    {\Huge\bfseries Python Programming\par}
    \vspace{0.8cm}
    {\LARGE Day 7: NumPy 기초\par}

    \vspace{2cm}

    {\large ndarray, Array Creation, Indexing, Vectorized Operations, Broadcasting, Linear Algebra\par}

    \vfill
\end{titlepage}

\tableofcontents
\newpage

% =================================================
\section{학습 목표}
% =================================================

\begin{learningbox}
이 장을 마치면 다음 내용을 설명하고 직접 코드로 작성할 수 있다.
\begin{itemize}[leftmargin=1.5em]
    \item NumPy의 목적과 \texttt{ndarray}의 의미 설명하기
    \item \texttt{np.array()}로 NumPy 배열 생성하기
    \item \texttt{zeros()}, \texttt{ones()}, \texttt{arange()}, \texttt{linspace()} 사용하기
    \item 1차원 배열에서 인덱싱과 슬라이싱 사용하기
    \item 배열 사이의 원소별 연산과 스칼라 연산 수행하기
    \item 브로드캐스팅을 이용하여 크기가 다른 배열 계산하기
    \item \texttt{@}, \texttt{.T}, \texttt{det()}, \texttt{inv()}로 기본 행렬 연산 수행하기
\end{itemize}
\end{learningbox}

% =================================================
\section{NumPy와 \texttt{ndarray}}
% =================================================

NumPy는 Numerical Python의 줄임말로, Python에서 수치 계산을 효율적으로 수행하기 위한 라이브러리이다. Python 리스트도 숫자를 저장할 수 있지만, NumPy 배열은 대량의 수치 데이터를 빠르게 처리하고 행렬 및 선형대수 계산을 수행하는 데 적합하다.

NumPy를 설치하지 않았다면 터미널에서 다음 명령을 한 번 실행한다.

\begin{verbatim}
pip install numpy
\end{verbatim}

NumPy는 일반적으로 \texttt{np}라는 별칭으로 불러온다.

\begin{lstlisting}
import numpy as np
\end{lstlisting}

\subsection{\texttt{ndarray}}

NumPy의 기본 배열 자료형은 \texttt{ndarray}이다. \texttt{np.array()}에 Python 리스트를 전달하면 NumPy 배열을 만들 수 있다.

\subsection{최종 코드: \texttt{01\_numpy\_intro.py}}

\begin{lstlisting}
import numpy as np

numbers = np.array([1, 2, 3, 4, 5])

print(numbers)
print(type(numbers))
\end{lstlisting}

실행 결과는 다음과 같다.

\begin{verbatim}
[1 2 3 4 5]
<class 'numpy.ndarray'>
\end{verbatim}

Python 리스트는 원소 사이에 쉼표가 표시되지만, NumPy 배열은 기본 출력에서 쉼표 없이 표시된다.

\begin{importantbox}
NumPy의 대부분 기능은 \texttt{ndarray}를 중심으로 동작한다. 이후의 인덱싱, 벡터화 연산, 브로드캐스팅, 행렬 연산도 모두 NumPy 배열을 대상으로 수행한다.
\end{importantbox}

% =================================================
\section{NumPy 배열 생성}
% =================================================

NumPy는 리스트를 변환하는 것뿐 아니라 일정한 규칙을 가진 배열을 빠르게 생성하는 함수도 제공한다.

\begin{center}
\begin{tabular}{ll}
\toprule
함수 & 기능 \\
\midrule
\texttt{np.zeros(n)} & 0으로 채운 길이 \texttt{n}의 배열 생성 \\
\texttt{np.ones(n)} & 1로 채운 길이 \texttt{n}의 배열 생성 \\
\texttt{np.arange(start, stop, step)} & 일정한 간격의 배열 생성 \\
\texttt{np.linspace(start, stop, count)} & 구간을 균등하게 나눈 배열 생성 \\
\bottomrule
\end{tabular}
\end{center}

\subsection{\texttt{zeros()}와 \texttt{ones()}}

\texttt{np.zeros(5)}는 0을 다섯 개 가진 배열을 만들고, \texttt{np.ones(5)}는 1을 다섯 개 가진 배열을 만든다. 기본 자료형은 실수이므로 \texttt{0.}, \texttt{1.}처럼 출력된다.

\subsection{\texttt{arange()}}

\texttt{np.arange(1, 11, 2)}는 1부터 시작하여 2씩 증가하고 11 직전에서 멈춘다. Python의 \texttt{range()}와 마찬가지로 끝값은 포함하지 않는다.

\subsection{\texttt{linspace()}}

\texttt{np.linspace(0, 1, 5)}는 0부터 1까지의 구간을 균등하게 나누어 총 다섯 개의 값을 만든다. 끝값도 포함된다는 점이 \texttt{arange()}와 다르다.

\subsection{최종 코드: \texttt{02\_numpy\_array.py}}

\begin{lstlisting}
import numpy as np

zeros = np.zeros(5)
ones = np.ones(5)
numbers = np.arange(1, 11, 2)
values = np.linspace(0, 1, 5)

print("Zeros :", zeros)
print("Ones :", ones)
print("Arange :", numbers)
print("Linspace :", values)
\end{lstlisting}

실행 결과는 다음과 같다.

\begin{verbatim}
Zeros : [0. 0. 0. 0. 0.]
Ones : [1. 1. 1. 1. 1.]
Arange : [1 3 5 7 9]
Linspace : [0.   0.25 0.5  0.75 1.  ]
\end{verbatim}

\begin{importantbox}
\texttt{arange()}는 간격을 기준으로 값을 만들고, \texttt{linspace()}는 전체 개수를 기준으로 값을 만든다. 그래프나 물리 시뮬레이션에서 일정 구간을 샘플링할 때는 \texttt{linspace()}가 자주 사용된다.
\end{importantbox}

% =================================================
\section{인덱싱과 슬라이싱}
% =================================================

NumPy 배열은 Python 리스트처럼 인덱싱과 슬라이싱을 지원한다. 인덱스는 0부터 시작하며 음수 인덱스로 뒤쪽 원소에 접근할 수 있다.

\begin{lstlisting}
numbers[index]
numbers[start:end]
\end{lstlisting}

슬라이싱에서 시작 위치는 포함하고 끝 위치는 포함하지 않는다.

\subsection{최종 코드: \texttt{03\_numpy\_indexing.py}}

\begin{lstlisting}
import numpy as np

numbers = np.array([10, 20, 30, 40, 50])

print(numbers[0])
print(numbers[2])
print(numbers[-1])

print(numbers[1:4])
print(numbers[:3])
print(numbers[2:])
\end{lstlisting}

실행 결과는 다음과 같다.

\begin{verbatim}
10
30
50
[20 30 40]
[10 20 30]
[30 40 50]
\end{verbatim}

\subsection{NumPy 슬라이싱의 View}

Python 리스트를 슬라이싱하면 일반적으로 새로운 리스트가 만들어진다. 반면 NumPy의 슬라이싱 결과는 원본 배열의 일부를 바라보는 View인 경우가 많다. 따라서 슬라이스를 수정하면 원본 배열도 함께 바뀔 수 있다.

\begin{lstlisting}
part = numbers[1:4]
part[0] = 999
\end{lstlisting}

원본과 완전히 독립된 배열이 필요할 때는 \texttt{copy()}를 사용한다.

\begin{lstlisting}
part = numbers[1:4].copy()
\end{lstlisting}

\begin{importantbox}
NumPy 슬라이싱은 메모리를 절약하기 위해 원본 데이터를 공유할 수 있다. 원본 변경을 원하지 않는다면 \texttt{copy()}를 명시적으로 사용한다.
\end{importantbox}

% =================================================
\section{벡터화 연산}
% =================================================

Python 리스트에서 \texttt{+}는 리스트를 연결하지만, NumPy 배열에서 산술 연산자는 같은 위치의 원소끼리 계산한다. 이러한 방식을 벡터화 연산(vectorized operation)이라고 한다.

\begin{verbatim}
[1 2 3] + [4 5 6] = [5 7 9]
\end{verbatim}

반복문을 직접 작성하지 않고 배열 전체를 한 번에 계산할 수 있으며, NumPy 내부의 최적화된 코드가 실행되므로 대량의 수치 계산에 효율적이다.

\subsection{최종 코드: \texttt{04\_numpy\_operations.py}}

\begin{lstlisting}
import numpy as np

a = np.array([1, 2, 3])
b = np.array([4, 5, 6])

print(a + b)
print(a - b)
print(a * b)
print(a / b)

print()

print(a + 10)
print(a * 5)
\end{lstlisting}

실행 결과는 다음과 같다.

\begin{verbatim}
[5 7 9]
[-3 -3 -3]
[ 4 10 18]
[0.25 0.4  0.5 ]

[11 12 13]
[ 5 10 15]
\end{verbatim}

\subsection{원소별 곱과 행렬곱}

\texttt{a * b}는 같은 위치의 원소를 곱하는 원소별 곱이다. 선형대수의 행렬곱은 \texttt{@} 연산자를 사용한다.

\begin{center}
\begin{tabular}{ll}
\toprule
연산 & 의미 \\
\midrule
\texttt{A * B} & 원소별 곱 \\
\texttt{A @ B} & 행렬곱 \\
\bottomrule
\end{tabular}
\end{center}

\begin{importantbox}
NumPy에서는 가능한 한 Python 반복문 대신 배열 전체에 연산을 적용한다. 코드가 짧아지고 대량의 데이터 처리도 더 효율적이다.
\end{importantbox}

% =================================================
\section{브로드캐스팅}
% =================================================

브로드캐스팅(broadcasting)은 크기가 다른 배열 사이의 연산을 NumPy가 자동으로 확장하여 처리하는 기능이다.

예를 들어 배열과 하나의 숫자를 더하면 그 숫자가 모든 원소에 적용된다.

\begin{lstlisting}
a = np.array([1, 2, 3])
print(a + 10)
\end{lstlisting}

개념적으로는 10이 \texttt{[10, 10, 10]}처럼 확장된 것으로 볼 수 있지만, NumPy는 실제로 같은 값을 반복 저장하지 않고 연산 규칙을 적용한다.

\subsection{행렬과 벡터의 브로드캐스팅}

2차원 행렬의 각 행 길이와 1차원 벡터의 길이가 같다면 벡터를 각 행에 적용할 수 있다.

\subsection{최종 코드: \texttt{05\_numpy\_broadcasting.py}}

\begin{lstlisting}
import numpy as np

a = np.array([1, 2, 3])
b = np.array([10])

print(a + b)

print()

matrix = np.array([[1, 2, 3],
                   [4, 5, 6]])

vector = np.array([10, 20, 30])

print(matrix + vector)
\end{lstlisting}

실행 결과는 다음과 같다.

\begin{verbatim}
[11 12 13]

[[11 22 33]
 [14 25 36]]
\end{verbatim}

행렬과 벡터의 계산은 개념적으로 다음과 같다.

\begin{verbatim}
[[1 2 3]     [[10 20 30]     [[11 22 33]
 [4 5 6]]  +  [10 20 30]]  =  [14 25 36]]
\end{verbatim}

\begin{importantbox}
브로드캐스팅은 배열의 크기가 다르더라도 일정한 규칙에 따라 호환될 때만 가능하다. 호환되지 않는 모양의 배열을 연산하면 오류가 발생한다.
\end{importantbox}

% =================================================
\section{NumPy 선형대수}
% =================================================

NumPy는 행렬곱, 전치행렬, 행렬식, 역행렬과 같은 기본 선형대수 연산을 지원한다.

\begin{center}
\begin{tabular}{ll}
\toprule
표현 & 기능 \\
\midrule
\texttt{A @ B} & 행렬곱 \\
\texttt{A.T} & 전치행렬 \\
\texttt{np.linalg.det(A)} & 행렬식 \\
\texttt{np.linalg.inv(A)} & 역행렬 \\
\bottomrule
\end{tabular}
\end{center}

\subsection{최종 코드: \texttt{06\_numpy\_linear\_algebra.py}}

\begin{lstlisting}
import numpy as np

A = np.array([[1, 2],
              [3, 4]])

B = np.array([[5, 6],
              [7, 8]])

print("Matrix Multiplication")
print(A @ B)

print()

print("Transpose")
print(A.T)

print()

print("Determinant")
print(np.linalg.det(A))

print()

print("Inverse")
print(np.linalg.inv(A))
\end{lstlisting}

실행 결과는 다음과 같다.

\begin{verbatim}
Matrix Multiplication
[[19 22]
 [43 50]]

Transpose
[[1 3]
 [2 4]]

Determinant
-2.0000000000000004

Inverse
[[-2.   1. ]
 [ 1.5 -0.5]]
\end{verbatim}

부동소수점 계산에서는 수학적으로 $-2$인 값이 \texttt{-2.0000000000000004}처럼 출력될 수 있다. 이는 이진 부동소수점 표현에서 발생하는 정상적인 수치 오차이다.

\subsection{연립방정식 풀이}

연립방정식 $A\mathbf{x}=\mathbf{b}$를 풀 때 역행렬을 직접 계산하기보다 \texttt{np.linalg.solve(A, b)}를 사용하는 편이 일반적으로 더 효율적이고 수치적으로 안정적이다.

\begin{lstlisting}
x = np.linalg.solve(A, b)
\end{lstlisting}

\begin{importantbox}
\texttt{np.linalg.inv()}는 역행렬이 존재하는 정사각행렬에만 사용할 수 있다. 실제 연립방정식 풀이에서는 역행렬을 직접 구하기보다 \texttt{solve()}를 우선 고려한다.
\end{importantbox}

% =================================================
\section{Day 7 종합 실습: NumPy 기초 복습}
% =================================================

이번 종합 실습에서는 배열 생성, 인덱싱, 벡터화 연산, 브로드캐스팅, 선형대수 기능을 한 파일에서 복습한다.

\subsection{프로그램 요구사항}

\begin{enumerate}[leftmargin=1.8em]
    \item \texttt{[10, 20, 30, 40, 50]} 배열과 자료형을 출력한다.
    \item \texttt{zeros()}, \texttt{ones()}, \texttt{arange()}, \texttt{linspace()}로 배열을 생성한다.
    \item 지정된 배열에서 첫 원소, 마지막 원소, 세 번째 원소와 여러 슬라이스를 출력한다.
    \item 두 배열의 덧셈, 뺄셈, 원소별 곱, 나눗셈과 스칼라 연산을 수행한다.
    \item 2차원 행렬과 1차원 벡터를 브로드캐스팅하여 더한다.
    \item 두 행렬의 행렬곱, 전치행렬, 행렬식, 역행렬을 출력한다.
\end{enumerate}

\subsection{최종 코드: \texttt{practice\_day07.py}}

\begin{lstlisting}
import numpy as np

# 1
v = np.array([10, 20, 30, 40, 50])
print(v)
print(type(v))

# 2
print(np.zeros(4))
print(np.ones(6))
print(np.arange(2, 12, 2))
print(np.linspace(0, 100, 5))

# 3
numbers = np.array([5, 10, 15, 20, 25])

print(numbers[0])
print(numbers[-1])
print(numbers[2])
print(numbers[:3])
print(numbers[-2:])
print(numbers[1:4])

# 4
a = np.array([2, 4, 6])
b = np.array([1, 3, 5])

print(a + b)
print(a - b)
print(a * b)
print(a / b)
print(a + 10)
print(b * 2)

# 5
matrix = np.array([
    [1, 2, 3],
    [4, 5, 6]
])

vector = np.array([10, 20, 30])

print(matrix + vector)

# 6
A = np.array([
    [2, 1],
    [5, 3]
])

B = np.array([
    [1, 2],
    [3, 4]
])

print(A @ B)
print(A.T)
print(np.linalg.det(A))
print(np.linalg.inv(A))
\end{lstlisting}

\subsection{실행 결과}

\begin{verbatim}
[10 20 30 40 50]
<class 'numpy.ndarray'>
[0. 0. 0. 0.]
[1. 1. 1. 1. 1. 1.]
[ 2  4  6  8 10]
[  0.  25.  50.  75. 100.]
5
25
15
[ 5 10 15]
[20 25]
[10 15 20]
[ 3  7 11]
[1 1 1]
[ 2 12 30]
[2.         1.33333333 1.2       ]
[12 14 16]
[ 2  6 10]
[[11 22 33]
 [14 25 36]]
[[ 5  8]
 [14 22]]
[[2 5]
 [1 3]]
1.0000000000000009
[[ 3. -1.]
 [-5.  2.]]
\end{verbatim}

\subsection{프로그램 구조 분석}

\begin{center}
\begin{tabular}{ll}
\toprule
기능 & 사용한 개념 \\
\midrule
기본 배열 생성 & \texttt{np.array()}와 \texttt{ndarray} \\
규칙적인 배열 생성 & \texttt{zeros()}, \texttt{ones()}, \texttt{arange()}, \texttt{linspace()} \\
일부 원소 선택 & 인덱싱과 슬라이싱 \\
배열 전체 계산 & 벡터화 연산 \\
행렬과 벡터 계산 & 브로드캐스팅 \\
행렬곱 및 역행렬 계산 & \texttt{np.linalg} \\
\bottomrule
\end{tabular}
\end{center}

\begin{importantbox}
NumPy의 핵심은 배열을 만든 뒤 반복문 없이 배열 전체에 연산을 적용하는 것이다. 필요한 고급 기능은 실제 프로젝트나 전공 문제에서 사용할 때 추가로 학습하면 된다.
\end{importantbox}

% =================================================
\section{실습 파일 목록}
% =================================================

\begin{practicebox}
Day 7의 학습 파일은 다음과 같이 관리한다.
\begin{itemize}[leftmargin=1.5em]
    \item \texttt{01\_numpy\_intro.py}
    \item \texttt{02\_numpy\_array.py}
    \item \texttt{03\_numpy\_indexing.py}
    \item \texttt{04\_numpy\_operations.py}
    \item \texttt{05\_numpy\_broadcasting.py}
    \item \texttt{06\_numpy\_linear\_algebra.py}
    \item \texttt{practice\_day07.py}
    \item \texttt{python\_day07.tex}
    \item \texttt{python\_day07.pdf}
\end{itemize}
\end{practicebox}

% =================================================
\section{핵심 요약}
% =================================================

\begin{itemize}[leftmargin=1.5em]
    \item NumPy는 Python에서 효율적인 수치 계산을 제공하는 라이브러리이다.
    \item NumPy의 기본 배열 자료형은 \texttt{ndarray}이다.
    \item \texttt{np.array()}는 Python 리스트를 NumPy 배열로 변환한다.
    \item \texttt{zeros()}와 \texttt{ones()}는 일정한 값으로 채운 배열을 생성한다.
    \item \texttt{arange()}는 간격을 기준으로, \texttt{linspace()}는 개수를 기준으로 배열을 생성한다.
    \item NumPy 배열은 Python 리스트처럼 인덱싱과 슬라이싱을 지원한다.
    \item NumPy 슬라이싱은 원본 데이터를 공유하는 View일 수 있다.
    \item 배열 연산은 같은 위치의 원소끼리 계산되는 벡터화 연산이다.
    \item 배열과 하나의 숫자를 연산하면 그 숫자가 모든 원소에 적용된다.
    \item 브로드캐스팅은 호환되는 서로 다른 크기의 배열을 자동으로 확장하여 계산한다.
    \item \texttt{*}는 원소별 곱이고 \texttt{@}는 행렬곱이다.
    \item \texttt{.T}는 전치행렬을 반환한다.
    \item \texttt{np.linalg.det()}는 행렬식, \texttt{np.linalg.inv()}는 역행렬을 계산한다.
    \item 실제 연립방정식 풀이에는 \texttt{np.linalg.solve()}가 더 적합한 경우가 많다.
    \item NumPy의 고급 기능은 프로젝트와 전공 문제에서 필요한 시점에 추가 학습한다.
\end{itemize}

\end{document}
